\iffalse
!!!!!!!!!!!!!!!!!!!!!!!!!!!!!!!!!!!!!
! Do not modify this file directly  !
! This file has been auto-generated !
!!!!!!!!!!!!!!!!!!!!!!!!!!!!!!!!!!!!!
\fi

%% ODER: format ==         = "\mathrel{==}"
%% ODER: format /=         = "\neq "
%
%
\makeatletter
\@ifundefined{lhs2tex.lhs2tex.sty.read}%
  {\@namedef{lhs2tex.lhs2tex.sty.read}{}%
   \newcommand\SkipToFmtEnd{}%
   \newcommand\EndFmtInput{}%
   \long\def\SkipToFmtEnd#1\EndFmtInput{}%
  }\SkipToFmtEnd

\newcommand\ReadOnlyOnce[1]{\@ifundefined{#1}{\@namedef{#1}{}}\SkipToFmtEnd}
\usepackage{newtxmath}
\usepackage{stmaryrd}
\DeclareFontFamily{OT1}{cmtex}{}
\DeclareFontShape{OT1}{cmtex}{m}{n}
  {<5><6><7><8>cmtex8
   <9>cmtex9
   <10><10.95><12><14.4><17.28><20.74><24.88>cmtex10}{}
\DeclareFontShape{OT1}{cmtex}{m}{it}
  {<-> ssub * cmtt/m/it}{}
\newcommand{\texfamily}{\fontfamily{cmtex}\selectfont}
\DeclareFontShape{OT1}{cmtt}{bx}{n}
  {<5><6><7><8>cmtt8
   <9>cmbtt9
   <10><10.95><12><14.4><17.28><20.74><24.88>cmbtt10}{}
\DeclareFontShape{OT1}{cmtex}{bx}{n}
  {<-> ssub * cmtt/bx/n}{}
\newcommand{\tex}[1]{\text{\texfamily#1}}	% NEU

\newcommand{\Sp}{\hskip.33334em\relax}


\newcommand{\Conid}[1]{\mathit{#1}}
\newcommand{\Varid}[1]{\mathit{#1}}
\newcommand{\anonymous}{\kern0.06em \vbox{\hrule\@width.5em}}
\newcommand{\plus}{\mathbin{+\!\!\!+}}
\newcommand{\bind}{\mathbin{>\!\!\!>\mkern-6.7mu=}}
\newcommand{\rbind}{\mathbin{=\mkern-6.7mu<\!\!\!<}}% suggested by Neil Mitchell
\newcommand{\sequ}{\mathbin{>\!\!\!>}}
\renewcommand{\leq}{\leqslant}
\renewcommand{\geq}{\geqslant}
\usepackage{polytable}

%mathindent has to be defined
\@ifundefined{mathindent}%
  {\newdimen\mathindent\mathindent\leftmargini}%
  {}%

\def\resethooks{%
  \global\let\SaveRestoreHook\empty
  \global\let\ColumnHook\empty}
\newcommand*{\savecolumns}[1][default]%
  {\g@addto@macro\SaveRestoreHook{\savecolumns[#1]}}
\newcommand*{\restorecolumns}[1][default]%
  {\g@addto@macro\SaveRestoreHook{\restorecolumns[#1]}}
\newcommand*{\aligncolumn}[2]%
  {\g@addto@macro\ColumnHook{\column{#1}{#2}}}

\resethooks

\newcommand{\onelinecommentchars}{\quad-{}- }
\newcommand{\commentbeginchars}{\enskip\{-}
\newcommand{\commentendchars}{-\}\enskip}

\newcommand{\visiblecomments}{%
  \let\onelinecomment=\onelinecommentchars
  \let\commentbegin=\commentbeginchars
  \let\commentend=\commentendchars}

\newcommand{\invisiblecomments}{%
  \let\onelinecomment=\empty
  \let\commentbegin=\empty
  \let\commentend=\empty}

\visiblecomments

\newlength{\blanklineskip}
\setlength{\blanklineskip}{0.66084ex}

\newcommand{\hsindent}[1]{\quad}% default is fixed indentation
\let\hspre\empty
\let\hspost\empty
\newcommand{\NB}{\textbf{NB}}
\newcommand{\Todo}[1]{$\langle$\textbf{To do:}~#1$\rangle$}

\EndFmtInput
\makeatother
%
%
%
%
%
%
% This package provides two environments suitable to take the place
% of hscode, called "plainhscode" and "arrayhscode". 
%
% The plain environment surrounds each code block by vertical space,
% and it uses \abovedisplayskip and \belowdisplayskip to get spacing
% similar to formulas. Note that if these dimensions are changed,
% the spacing around displayed math formulas changes as well.
% All code is indented using \leftskip.
%
% Changed 19.08.2004 to reflect changes in colorcode. Should work with
% CodeGroup.sty.
%
\ReadOnlyOnce{polycode.fmt}%
\makeatletter

\newcommand{\hsnewpar}[1]%
  {{\parskip=0pt\parindent=0pt\par\vskip #1\noindent}}

% can be used, for instance, to redefine the code size, by setting the
% command to \small or something alike
\newcommand{\hscodestyle}{}

% The command \sethscode can be used to switch the code formatting
% behaviour by mapping the hscode environment in the subst directive
% to a new LaTeX environment.

\newcommand{\sethscode}[1]%
  {\expandafter\let\expandafter\hscode\csname #1\endcsname
   \expandafter\let\expandafter\endhscode\csname end#1\endcsname}

% "compatibility" mode restores the non-polycode.fmt layout.

\newenvironment{compathscode}%
  {\par\noindent
   \advance\leftskip\mathindent
   \hscodestyle
   \let\\=\@normalcr
   \let\hspre\(\let\hspost\)%
   \pboxed}%
  {\endpboxed\)%
   \par\noindent
   \ignorespacesafterend}

\newcommand{\compaths}{\sethscode{compathscode}}

% "plain" mode is the proposed default.
% It should now work with \centering.
% This required some changes. The old version
% is still available for reference as oldplainhscode.

\newenvironment{plainhscode}%
  {\hsnewpar\abovedisplayskip
   \advance\leftskip\mathindent
   \hscodestyle
   \let\hspre\(\let\hspost\)%
   \pboxed}%
  {\endpboxed%
   \hsnewpar\belowdisplayskip
   \ignorespacesafterend}

\newenvironment{oldplainhscode}%
  {\hsnewpar\abovedisplayskip
   \advance\leftskip\mathindent
   \hscodestyle
   \let\\=\@normalcr
   \(\pboxed}%
  {\endpboxed\)%
   \hsnewpar\belowdisplayskip
   \ignorespacesafterend}

% Here, we make plainhscode the default environment.

\newcommand{\plainhs}{\sethscode{plainhscode}}
\newcommand{\oldplainhs}{\sethscode{oldplainhscode}}
\plainhs

% The arrayhscode is like plain, but makes use of polytable's
% parray environment which disallows page breaks in code blocks.

\newenvironment{arrayhscode}%
  {\hsnewpar\abovedisplayskip
   \advance\leftskip\mathindent
   \hscodestyle
   \let\\=\@normalcr
   \(\parray}%
  {\endparray\)%
   \hsnewpar\belowdisplayskip
   \ignorespacesafterend}

\newcommand{\arrayhs}{\sethscode{arrayhscode}}

% The mathhscode environment also makes use of polytable's parray 
% environment. It is supposed to be used only inside math mode 
% (I used it to typeset the type rules in my thesis).

\newenvironment{mathhscode}%
  {\parray}{\endparray}

\newcommand{\mathhs}{\sethscode{mathhscode}}

% texths is similar to mathhs, but works in text mode.

\newenvironment{texthscode}%
  {\(\parray}{\endparray\)}

\newcommand{\texths}{\sethscode{texthscode}}

% The framed environment places code in a framed box.

\def\codeframewidth{\arrayrulewidth}
\RequirePackage{calc}

\newenvironment{framedhscode}%
  {\parskip=\abovedisplayskip\par\noindent
   \hscodestyle
   \arrayrulewidth=\codeframewidth
   \tabular{@{}|p{\linewidth-2\arraycolsep-2\arrayrulewidth-2pt}|@{}}%
   \hline\framedhslinecorrect\\{-1.5ex}%
   \let\endoflinesave=\\
   \let\\=\@normalcr
   \(\pboxed}%
  {\endpboxed\)%
   \framedhslinecorrect\endoflinesave{.5ex}\hline
   \endtabular
   \parskip=\belowdisplayskip\par\noindent
   \ignorespacesafterend}

\newcommand{\framedhslinecorrect}[2]%
  {#1[#2]}

\newcommand{\framedhs}{\sethscode{framedhscode}}

% The inlinehscode environment is an experimental environment
% that can be used to typeset displayed code inline.

\newenvironment{inlinehscode}%
  {\(\def\column##1##2{}%
   \let\>\undefined\let\<\undefined\let\\\undefined
   \newcommand\>[1][]{}\newcommand\<[1][]{}\newcommand\\[1][]{}%
   \def\fromto##1##2##3{##3}%
   \def\nextline{}}{\) }%

\newcommand{\inlinehs}{\sethscode{inlinehscode}}

% The joincode environment is a separate environment that
% can be used to surround and thereby connect multiple code
% blocks.

\newenvironment{joincode}%
  {\let\orighscode=\hscode
   \let\origendhscode=\endhscode
   \def\endhscode{\def\hscode{\endgroup\def\@currenvir{hscode}\\}\begingroup}
   %\let\SaveRestoreHook=\empty
   %\let\ColumnHook=\empty
   %\let\resethooks=\empty
   \orighscode\def\hscode{\endgroup\def\@currenvir{hscode}}}%
  {\origendhscode
   \global\let\hscode=\orighscode
   \global\let\endhscode=\origendhscode}%

\makeatother
\EndFmtInput
%
%


%
%
%
%
%
\ReadOnlyOnce{lineno.fmt}%

\usepackage{etoolbox}

\makeatletter

\newtoggle{lineno@column}
\newtoggle{lineno@print}
\newcounter{lineno@num}
\newlength\linenowidth

\newcommand\linenoFormat{\footnotesize}

\newcommand\linenoSetup{%
  \iftoggle{lineno@column}{%
    \column{lineno}{@{}r}%
    \column{B}{}
  }{}%
}

\newcommand\linenoStart{%
  \refstepcounter{lineno@num}%
  \iftoggle{lineno@print}{%
    \fromto{lineno}{B}{\makebox[\linenowidth][r]{\linenoFormat\arabic{lineno@num}}}%
  }{}%
}

\newcommand\linenoon[1][]{%
  #1\toggletrue{lineno@column}%
  #1\toggletrue{lineno@print}%
}

\newcommand\linenooff[1][]{%
  #1\togglefalse{lineno@column}%
  #1\togglefalse{lineno@print}%
}

\newcommand\linenoskip[1][]{%
  #1\toggletrue{lineno@column}%
  #1\togglefalse{lineno@print}%
}

\newcommand\linenoreset[1][1]{%
  \defcounter{lineno@num}{#1 - 1}%
}

\EndFmtInput












% Some freest specific formatting directives

\appendix


The organization of the supplement follows the structure of the main document.
Topics are covered in the same sequence, so the two documents can be read
side-by-side.


\section{Motivation --- Remarks}
\label{sec:remarks}

%\subsection{Remarks}

% Switch to the borrow code style for this appendix section
{\BorrowCodeStyle

The reader may ask why \ensuremath{\CodeOp{select}} still adheres to
resource-passing style. The answer is simplicity of our system. While it
is possible to define \ensuremath{\CodeOp{select}} with a borrowing-based
interface, supporting it requires subtyping of session types. While
subtyping is fine in theory, it turns out that it is undecidable for
CFST~\cite{DBLP:journals/toplas/Padovani19}. Hence, it cannot be the
basis for a complete type checker.

For borrowing from regular session
types~\cite{DBLP:journals/pacmpl/SaffrichS0V25}, a borrowing-style
\ensuremath{\CodeOp{select}} operation is possible because subtyping for regular
session types is decidable.

A borrowing-style \ensuremath{\CodeOp{branch}} operator (external choice) cannot be supported in \thecalculus.
This operator selects between different continuations at run time and
the residual channel depends on this run-time choice.

% The downside is that some uses of |select| force us to use
% the |drop| primitive, which discharges an exhausted channel
% borrow of type |Drop|. The discharge is part of the
% |send| and |recv| operations, so it is rarely
% needed in practice. Alternatively, the type checker could insert
% |drop| automatically.

All \thecalculus examples are available as part of our supplement. Compared to the
presentation in the main text, the split operators carry type
parameters in the actual code to simplify type checking.

% End the borrow code style for this appendix section.
}

\clearpage

\section{Types: Session-Type Laws and Well Formedness Predicates}
\label{sec:appendix-types}

\begin{figure}[h]
  \begin{mathpar}
    \RuleUUnit \and
    \RuleUArr \and
    \RuleUProd \and
    \RuleUVariant \\
    \RuleUCtxEmpty \and
    \RuleUCtxVar \and
    \RuleUCtxSeq \and
    \RuleUCtxPar \\
    \RuleSNewSkip \and
    \RuleSNewSend \and
    \RuleSNewRecv \and
    \RuleSNewVar \\
    \RuleSNewSeq \and
    \RuleSNewBranch \and
    \RuleSNewChoice \and
    \RuleSNewRec
  \end{mathpar}
  \caption{Unrestricted types and contexts ($\Unr\T$, $\Unr\Ctx$); session
    types valid for $\CNew$ ($\SNew\S$)}
  \label{fig:omitted-predicates}
\end{figure}

\vfill

\begin{figure}[h]
  \medskip
  Laws for context-free session types.
  \begin{mathpar}
    \SSeq\SSkip\S = \S \and
    \SSeq \S\SSkip = \S \and
    \SSeq{\S[1]}{(\SSeq{\S[2]}{\S[3]})} =
    \SSeq{(\SSeq{\S[1]}{\S[2]})}{{\S[3]}} \and
    \SRec\SVar\S = \S{[\SRec\SVar\S/\SVar]} \and
    \SSeq{(\SChoice{\ell:\S[\ell]}{\ell\in L})}\S =
    \SChoice{\ell:\SSeq{\S[\ell]}\S}{\ell\in L} \and
    \SSeq{(\SBranch{\ell:\S[\ell]}{\ell\in L})}\S = \SBranch{\ell:\SSeq{\S[\ell]}\S}{\ell\in L}
  \end{mathpar}

  \medskip
  Laws for duality $\SDual\S$ (there is no dual for $\SDrop$ and $\SAcq$).
  \begin{mathpar}
    \SDual\SSkip = \SSkip \and
    \SDual{(\SSeq{\S[1]}{\S[2]})} = \SSeq{\SDual{\S[1]}}{\SDual{\S[2]}}
    \and
    \SDual{\SSend\T} = \SRecv\T \and
    \SDual{\SRecv\T} = \SSend\T \and
    \SDual{(\SChoice{\ell:\S[\ell]}{\ell\in L})} =
    \SBranch{\ell:\SDual{\S[\ell]}}{\ell\in L} \and
    \SDual{(\SBranch{\ell:\S[\ell]}{\ell\in L})} =
    \SChoice{\ell:\SDual{\S[\ell]}}{\ell\in L} \and
    \SDual\STerm = \SWait \and
    \SDual\SWait = \STerm \and
    \SDual{(\SRec\SVar\S)} = \SDual{\S{[\SRec\SVar\S/\SVar]}}
  \end{mathpar}
  \caption{Session-type equality and duality}
  \label{fig:session-laws}
\end{figure}

\vfill\vfill
\clearpage

\section{Dynamics: Congruence and Reduction Rules}

\begin{figure}[h]
\begin{align*}
  \U,\V \grmeq
    & \C
      \grmor \EVar
      \grmor \EAbs \EVar\E
      \grmor \EInj{i} \V
      \grmor \EPair \U\V
      % \grmor \CSend \V 
              \tag{Values}
  \\
  \EC \grmeq&
              \ECEmpty
              \grmor \EAppD[\DirLeft] \EC\E
              \grmor \EAppD[\DirLeft] \V\EC
              \grmor \EAppD[\DirRight] \E\EC
              \grmor \EAppD[\DirRight] \EC\V
              \grmor \ESeq \EC\E
              \grmor \EPair \EC\E
              \grmor \EPair \V\EC
              \\ \grmor
    & \EInj{i}\EC \grmor \ELet {\EPair \EVar\EVar}\EC\E
              \grmor \ECase \EC {\ECaseBranch* i \EVar \E}
              \tag{Evaluation contexts}
  \\
  \FC \grmeq
  & \PExp \EC
    \tag{Process evaluation contexts}
  \\
  \CC \grmeq
  & \ECEmpty
    \grmor \PPar{\CC}{\Proc}
    \grmor \PNew \BindGroup\BindGroup \CC
    \tag{Process contexts}
\end{align*}
  \caption{Values and evaluation contexts}
  \label{fig:values-evaluation-contexts}
\end{figure}

\vfill

\begin{figure}[h]
  \begin{mathpar}
    \RuleExpRedApp \and
    \RuleExpRedSeq \and
    \RuleExpRedPairElim \and
    \RuleExpRedUnfold \and
    \RuleExpRedVariantElim \and
    \RuleExpRedCtx
  \end{mathpar}
  \caption{Expression reduction ($\E \EReduce \E$)}
  \label{fig:expression-reduction}
\end{figure}

\vfill

\begin{figure}[h]
  \begin{mathpar}
    % \RuleProcCongRefl \and
    % \RuleProcCongTrans \and
    \RuleProcCongParComm \and
    \RuleProcCongParAssoc \and
    \RuleProcCongParUnit \and
    \RuleProcCongResSwap \and
    \RuleProcCongResComm \and
    \RuleProcCongExtend
  \end{mathpar}
  \caption{Process congruence}
  \label{fig:configuration-congruence}
\end{figure}

\vfill
\clearpage


\section{Translation: Congruence of Translated Processes}

\begin{figure}[h]
  \begin{mathpar}
    %\RuleTransCongRefl \and
    \RuleTransCongParComm \and
    \RuleTransCongParAssoc \and
    \RuleTransCongParUnit \and
    \RuleTransCongResSwap \and
    \RuleTransCongResComm \and
    \RuleTransCongFlagComm \and
    \RuleTransCongResFlagComm \and
    \RuleTransCongResExtend \and
    \RuleTransCongFlagExtend
    %\RuleTransCongTrans
  \end{mathpar}
  \caption{Congruence rules of translated processes}
  \label{fig:translated-congruence}
\end{figure}

\clearpage


\section{Algorithmic Typing}
\label{sec:algorithmic-typing-appendix}

\begin{figure}[h]
  \begin{minipage}{0.49\linewidth}
  \[
    \Ctx[1] \CtxJoin \Ctx[2] = \begin{cases}
      \Ctx[1] \CtxPar \Ctx[2] & \text{if $\Dir \in \{ \DirUnord, \DirNone \}$} \\
      \Ctx[1] \CtxSeq \Ctx[2] & \text{if $\Dir = \DirLeft$} \\
      \Ctx[2] \CtxSeq \Ctx[1] & \text{if $\Dir = \DirRight$}
    \end{cases}
  \]
  \end{minipage}
  \begin{minipage}{0.49\linewidth}
  \[
    \CtxRestrict*{X} = \begin{cases}
      \CtxRestrict*[1]{X} \CtxPar \CtxRestrict*[2]{X} & \text{if $\Ctx=\Ctx[1] \CtxPar \Ctx[2]$} \\
      \CtxRestrict*[1]{X} \CtxSeq \CtxRestrict*[2]{X} & \text{if $\Ctx=\Ctx[1] \CtxSeq \Ctx[2]$} \\
      \CtxVar \EVar \T & \text{if $\Ctx=\CtxVar \EVar T$ and $\EVar \in X$} \\
      \CtxEmpty & \text{otherwise}
    \end{cases}
  \]
  \end{minipage}
  \[
    \CaseJoinDir(\Ctx, X) = \begin{cases}
      \DirUnord & \text{if $\CtxRestrict*{X} \CtxPar \CtxRestrict*{X^C} \SubCtx \Ctx$}\\
      \DirLeft & \text{if $\CtxRestrict*{X} \CtxSeq \CtxRestrict*{X^C} \SubCtx \Ctx$}\\
      \text{undef.} & \text{otherwise}\\
    \end{cases}
  \]
  \caption{Context join $\Ctx[1] \CtxJoin \Ctx[2]$ joins two contexts
    according to the direction $\Dir$, context restriction
    $\CtxRestrict*{X}$ restricts a context to a set of variables $X$, and $\CaseJoinDir$ calculates if a context $\Ctx$ restricted to a set $X$ can be joined parallelly with $\Ctx$ restricted to the complement set.}
  \label{fig:context-operators}
\end{figure}

\vfill

\begin{figure}[h]
  \begin{mathpar}
    \RuleCSPair \and
    \RuleCSSum \and
    \RuleCSUnit \and
    \RuleCSArrow \and
    \RuleCSSwap \and
    \RuleCSContext
  \end{mathpar}
  \caption{Constraint simplification}
  \label{fig:constraint-simplification}
\end{figure}

\vfill

\begin{figure}[h]
  \small
  \hfill \fbox{$\ConcatType {\T[in]} {\T[in]} = \T[out]$}
  \begin{mathpar}
    \TypeConcatenationSkip

    \TypeConcatenationSemiDist

    \TypeConcatenationNonSemi
  \end{mathpar}

  \hfill \fbox{$\NormaliseType {\T[in]} {\T[out]}$}
  \begin{mathpar}
    \TypeNormalisationSemiSkips
    
    \TypeNormalisationSemi

    \TypeNormalisationMu

    \TypeNormalisationChoiceSemi

    \TypeNormalisationBranchSemi

    \TypeNormalisationIdentity
  \end{mathpar}
  \caption{Type Normalisation}
  \label{fig:typing-normalisation}
\end{figure}

\vfill


\clearpage
\section{Implementation}
\label{sec:appendix-implementation}

\subsection{Constraint Checking with FreeST}

In the implementation, we use FreeST
\cite{DBLP:journals/corr/abs-1904-01284,DBLP:journals/iandc/AlmeidaMTV22,
  almeida26:_frees_concur_progr_languag_with} version 5.0 to perform equivalence
checking. To this end, we generate little FreeST programs that typecheck if and
only if the equivalence constraints hold. 


For every equivalence constraint $\T[1] \CEq \T[2]$, we generate the following
FreeST program consisting of type definitions for $\T[1]$ and $\T[2]$ as well as
an identity function mapping $\T[1]$ to $\T[2]$.
For example, let $\T[1] = \SRec{\SVar}{\SSeq {\SSeq {\SSend{\BInt}} {\SSend {\BString}}} {\SVar}}$ and $\T[2] = \SRec{\SVar}{\SSeq {\SSkip} {\SSeq {\SSeq {\SSend{\BInt}} {\SSend {\BString}}} {\SVar}}}$. For the constraint $\T[1] \CEq \T[2]$ we generate the following FreeST program:

{
  \setlength{\blanklineskip}{\medskipamount}
  \TightCode[\bigskipamount]
\begin{hscode}\SaveRestoreHook\linenoSetup\linenoStart
\column{B}{@{}>{\hspre}c<{\hspost}@{}}%
\column{BE}{@{}l@{}}%
\column{7}{@{}>{\hspre}l<{\hspost}@{}}%
\column{E}{@{}>{\hspre}l<{\hspost}@{}}%
\>[B]{}\qquad {}\<[BE]%
\>[7]{}\CodeKw{type}\;\Conid{X\char95 T1}\mathbin{:}\mathrm{1S}{}\<[E]%
\\
\linenoStart \>[7]{}\CodeKw{type}\;\Conid{X\char95 T1}\mathrel{=}\mathord{!}\Conid{Int}\mathop{;}\mathord{!}\Conid{String}\mathop{;}\Conid{X\char95 T1}{}\<[E]%
\\[\blanklineskip]%
\linenoStart \>[7]{}\CodeKw{type}\;\Conid{T1}\mathbin{:}\mathrm{1S}{}\<[E]%
\\
\linenoStart \>[7]{}\CodeKw{type}\;\Conid{T1}\mathrel{=}\Conid{X\char95 T1}{}\<[E]%
\\[\blanklineskip]%
\linenoStart \>[7]{}\CodeKw{type}\;\Conid{X\char95 T2}\mathbin{:}\mathrm{1S}{}\<[E]%
\\
\linenoStart \>[7]{}\CodeKw{type}\;\Conid{X\char95 T2}\mathrel{=}\CodeTy{Skip}\mathop{;}\mathord{!}\Conid{Int}\mathop{;}\mathord{!}\Conid{String}\mathop{;}\Conid{X\char95 T2}{}\<[E]%
\\[\blanklineskip]%
\linenoStart \>[7]{}\CodeKw{type}\;\Conid{T2}\mathbin{:}\mathrm{1S}{}\<[E]%
\\
\linenoStart \>[7]{}\CodeKw{type}\;\Conid{T2}\mathrel{=}\Conid{X\char95 T2}{}\<[E]%
\\[\blanklineskip]%
\linenoStart \>[7]{}\Varid{id}\mathbin{:}\Conid{T1}\to \Conid{T2}{}\<[E]%
\\
\linenoStart \>[7]{}\Varid{id}\;\Varid{x}\mathrel{=}\Varid{x}{}\<[E]%
\ColumnHook
\end{hscode}\resethooks
}

\noindent FreeST 5.0 requires specifying kinds for type definitions.
Kind \ensuremath{\mathrm{1S}} says that the type needs to be a session type.
If \ensuremath{\Conid{T1}} and \ensuremath{\Conid{T2}} are the translated types from the equivalence constraint,
it is easy to see that this program only typechecks, if \ensuremath{\Conid{T1}} and
\ensuremath{\Conid{T2}} are equivalent session types.

The main parts of the translation are:
\begin{itemize}
  \item Translating inline recursion of the form $\SRec\SVar\S$ to explicit recursion.
  \item Translating mobility and direction annotations on function and pair
    types.
  \item Translating effects.
\end{itemize}

To translate a $\mu$-type to explicit recursion, we generate a fresh name and
associate the name with the recursion unrolled and translation applied.
We return the associated name, instead of the type.
We encode mobility, direction, and effects in a branch type with labels
indicating the state of the annotation.
Thus, the translation function $\tofreest{\T}$, defined in
\cref{fig:freest-translation}, takes a type $\T$ and returns a pair
$\langle F,U\rangle$ consisting of a FreeST type $F$ and a list of explicit
type definitions $U$.

\begin{figure}[tp]
  \footnotesize
  \setlength{\jot}{0.1ex}
  \begin{align*}
    \tofreest{\TypeKeyword{Unit}}
      & = \langle (), \emptyset \rangle \\
    \tofreest{\SSkip}
      & = \langle \SSkip, [] \rangle \\
    \tofreest{\STerm}
      & = \langle \STerm, [] \rangle \\
    \tofreest{\SWait}
      & = \langle \SWait, [] \rangle \\
    \tofreest{\TArr{\Mob}{\Dir, {\Eff}}{\T[1]}{\T[2]}}
      & = \langle (F_1 \to F_2) \TPair \SBranch {\ell : \SSkip} {\ell \in \text{dir} \cup \text{mob} \cup \text{eff}},\,
        U_1 \cup U_2 \rangle \quad \text{where} \\
      & \quad \langle F_1, U_1 \rangle = \tofreest{\T[1]} \\
      & \quad \langle F_2, U_2 \rangle = \tofreest{\T[2]} \\
      & \text{dir} = \begin{cases}
        \{ \text{left} \} & \text{if } \Dir = \DirLeft \\
        \{ \text{right} \} & \text{if } \Dir = \DirRight \\
        \{ \text{lin} \} & \text{if } \Dir = \DirUnord \\
        \emptyset & \text{if } \Dir = \DirNone \\
      \end{cases} \\
      & \text{mob} = \begin{cases}
        \{ \text{mobile} \} & \text{if } \Mob = \MobMobile \\
        \{ \text{static} \} & \text{if } \Mob = \MobStationary \\
      \end{cases} \\
      & \text{eff} = \begin{cases}
        \{ \text{impure} \} & \text{if } \Eff = \EffImpure \\
        \{ \text{} \} & \text{if } \Eff = \EffImpure \\
      \end{cases} \\
    \tofreest{\T[1] \TPair[\Dir] \T[2]}
      & = \langle (F_1 \times F_2) \TPair \SBranch {\ell : \SSkip} {\ell \in \text{dir}},\,
        U_1 \cup U_2 \rangle \quad \text{where} \\
      & \quad \langle F_1, U_1 \rangle = \tofreest{\T[1]} \\
      & \quad \langle F_2, U_2 \rangle = \tofreest{\T[2]} \\
      & \text{dir} = \begin{cases}
        \{ \text{left} \} & \text{if } \Dir = \DirLeft \\
        \{ \text{right} \} & \text{if } \Dir = \DirRight \\
        \{ \text{lin} \} & \text{if } \Dir = \DirUnord \\
        \emptyset & \text{if } \Dir = \DirNone
      \end{cases} \\
    \tofreest{\TVariant{\ell:\T[\ell]}}
      & = \langle \TVariant{\ell:F_\ell},\, \bigcup_{\ell\in L} U_\ell \rangle \quad \text{where} \\
      & \quad \langle F_\ell, U_\ell \rangle = \tofreest{\T[\ell]} \quad \text{for every } \ell \in L \\
    \tofreest{\SSend\T}
      & = \langle \SSend{F}, U \rangle \quad \text{where} \\
      & \quad \langle F, U \rangle = \tofreest{\T} \\
    \tofreest{\SRecv\T}
      & = \langle \SRecv{F}, U \rangle \quad \text{where} \\
      & \quad \langle F, U \rangle = \tofreest{\T} \\
    \tofreest{\S[1] ; \S[2]}
      & = \langle F_1 ; F_2,\, U_1 \cup U_2 \rangle \quad \text{where} \\
      & \quad \langle F_1, U_1 \rangle = \tofreest{\S[1]} \\
      & \quad \langle F_2, U_2 \rangle = \tofreest{\S[2]} \\
    \tofreest{\SChoice{\ell:\S[\ell]}{\ell\in L}}
      & = \langle +\{\ell : F_\ell\}_{\ell\in L},\, \bigcup_{\ell\in L} U_\ell \rangle \quad \text{where} \\
      & \quad \langle F_\ell, U_\ell \rangle = \tofreest{\S[\ell]} \quad \text{for every } \ell \in L \\
    \tofreest{\SBranch{\ell:\S[\ell]}{\ell\in L}}
      & = \langle \&\{\ell : F_\ell\}_{\ell\in L},\, \bigcup_{\ell\in L} U_\ell \rangle \quad \text{where} \\
      & \quad \langle F_\ell, U_\ell \rangle = \tofreest{\S[\ell]} \quad \text{for every } \ell \in L \\
    \tofreest{\SDrop}
      & = \langle \SSend{\TypeKeyword{Drop}},\, \emptyset \rangle \\
    \tofreest{\SAcq}
      & = \langle \SRecv{\TypeKeyword{Drop}}, \emptyset \rangle \\
    \tofreest{\SRec\SVar\S}
      & = \langle \SVar',\, U \cup \{\SVar' \mapsto U\} \rangle \quad \text{where } \SVar' \text{ is fresh and} \\
      & \quad \langle F, U \rangle = \tofreest{\S\,[\SVar'/\SVar]} \\
    \tofreest{\SVar}
      & = \langle \SVar, \emptyset \rangle
  \end{align*}
  \caption{Translation of types to FreeST types.}
  \label{fig:freest-translation}
\end{figure}

\subsection{Soundness and the Completeness of the translation}

\begin{definition}
  The equivalence relation $\simeq$ on \thecalculus types is the least
  congruence closed under the equivalence rules from
  \citet{DBLP:conf/icfp/ThiemannV16} extended with rules from
  \Cref{fig:cfstb-equivalence}. 
\end{definition}

\begin{figure}[tp]
  \begin{mathpar}
    \inferrule{
      \T[1] \simeq \T*[1] \\
      \T[2] \simeq \T*[2]
    }{
      \TArr{\Mob}{\Dir,\Eff}{\T[1]}{\T[2]}
      \simeq
      \TArr{\Mob}{\Dir,\Eff}{\T*[1]}{\T*[2]}
    }
    \and
    \inferrule{
      \T[1] \simeq \T*[1] \\
      \T[2] \simeq \T*[2]
    }{
      \T[1] \TPair \T[2]
      \simeq
      \T*[1] \TPair \T*[2]
    }
    \and
    \inferrule{
      \T[1] \simeq \T*[1] \\
      \T[2] \simeq \T*[2]
    }{
      \TSum{\T[1]}{\T[2]}
      \simeq
      \TSum{\T*[1]}{\T*[2]}
    }
  \end{mathpar}
  \caption{Modified rules for $\simeq$}
  \label{fig:cfstb-equivalence}
\end{figure}

% \begin{definition}[Substituting definitions]
%   Let $F$ be a FreeST session type and $U$ be a mapping from variables to FreeST
%   session types such that $\text{freevars}(F) \subseteq \text{dom}(U)$. We write
%   $F[U]$ for the FreeST type obtained by substituting the free variables in $F$
%   according to mapping $U$.
%   A well formed substitution has no free variables. 
% \end{definition}

Let $F$ be a FreeST session type and $U$ be a mapping from variables to FreeST
session types. In the following, we write $F[U]$ for the pair of $F$ and $U$
where $\text{freevars}(F) \subseteq \text{dom}(U)$. This notation includes
unfolding of definitions in the following way:
\begin{gather*}
  \frac{(X \mapsto F) \in U}{X[U] \simeq F[U]}
\end{gather*}

\begin{lemma}[Soundness of the translation]
  If $\T[1] \simeq \T[2]$, then $F_1 [U_1] \simeq F_2 [U_2]$
  where $\langle F_1, U_1 \rangle = \tofreest{\T[1]}$ and
  $\langle F_2, U_2 \rangle = \tofreest{\T[2]}$.
  \label{lem:sound-translation}
\end{lemma}
\begin{proof}
  We sketch a proof. The soundness proof requires two things: the translation of
  annotations in the arrow and product types preserves equivalence and the
  equivalence of the explicit recursion with the $\mu$-binder syntax. For the
  former we can perform induction on the equivalence of $\T[1] \simeq \T[2]$ and
  use equivalence and bisimulation lemmas from \citet{DBLP:conf/icfp/ThiemannV16}
  for regular type and session type cases, respectively. The pairs we introduce
  for annotations can be handled trivially by the type equivalence rules. To
  show explicit recursion translation preserves equivalence, we can use the laws
  for $\mu$-types from \citet{DBLP:conf/icfp/ThiemannV16}. 
\end{proof}

\begin{lemma}[Completeness of the translation]
  Let $\langle F_1, U_1 \rangle = \tofreest{\T[1]}$ and
  $\langle F_2, U_2 \rangle = \tofreest{\T[2]}$. If
  $F_1 [U_1] \simeq F_2 [U_2]$, then $\T[1] \simeq \T[2]$.
\end{lemma}
\begin{proof}
  Similar to \ref{lem:sound-translation} but the opposite direction.
\end{proof}

\clearpage

\section{Metatheory}
\label{sec:appendix-metatheory}

\begin{figure}[h]
  \begin{mathpar}
    \RuleExpValueBlocked \and
    \RuleExpConstBlocked \and
    % \RuleParReady \and
    % \RuleNuReady \and
    \RuleParBlocked \and
    \RuleNuBlocked \and
    \RuleNuBlockedAcq
    % \RuleNuBlockedAcq \and
    % \RuleNuBlockedThere
  \end{mathpar}
  where
  \[
    \begin{array}{ll}
      \BlockedChannelsSend & \BlockedChannelsRecv \\
      \BlockedChannelsSelect & \BlockedChannelsBranch \\
      \BlockedChannelsTerm & \BlockedChannelsWait \\
      \BlockedChannelsAcquire & \BlockedChannelsDrop \\
      \BlockedChannelsLSplit & \BlockedChannelsRSplit \\
      \BlockedChannelsFork & \BlockedChannelsNew \\
      \BlockedChannelsDiscard \\\\
      \AcqBlockedChannelsSend & \AcqBlockedChannelsRecv \\
      \AcqBlockedChannelsSelect & \AcqBlockedChannelsBranch \\
      \AcqBlockedChannelsTerm & \AcqBlockedChannelsWait \\
      \AcqBlockedChannelsAcquire & \AcqBlockedChannelsDrop \\
      \AcqBlockedChannelsLSplit & \AcqBlockedChannelsRSplit \\
      \AcqBlockedChannelsFork & \AcqBlockedChannelsNew \\
      \AcqBlockedChannelsDiscard 
    \end{array}
  \]
  \caption{Blocked Processes (Blocked $\Proc$)}
  \label{fig:blocked-expr-proc}
\end{figure}

\clearpage
\subsection{Type Safety for Declarative Typing}
\label{sec:appendix-metatheory-ts}

\providecommand{\proofcase}[1]{\smallskip\noindent\textit{Case}~#1\quad}

\begin{lemma}[Process progress]
  \label{lem:process-progress}
  If $\HasType[\EffImpure]\SessCtx\E\BUnit$, then $\E$ is a value, or
  $\E \EReduce \E*$ for some $\E*$, or $\E$ is blocked. If $\E$ is a value, then
  $\E = \CUnit$.
\end{lemma}
\begin{proof}
  The trichotomy is expression progress; a stuck redex $\ECapp{\EApp\C\V}$ is a
  blocked process. A value of type $\BUnit$ in a session context is $\CUnit$ by
  canonical forms.
\end{proof}

\begin{theorem}[Process progress, restated]
  \label{thm:process-progress-app}
  If $\ProcHasType\SessCtx\Proc$, then $\Proc \Congruent \PExp\CUnit$, or
  $\Blocked\Proc$, or $\Proc \PReduce \Proc*$ for some $\Proc*$.
\end{theorem}
\begin{proof}
  Induction on the derivation of $\ProcHasType\SessCtx\Proc$; \textsc{TP-Weaken}
  follows from the induction hypothesis.

  \proofcase{\textsc{TP-Expr}, $\Proc = \PExp\E$, $\HasType[\EffImpure]\SessCtx\E\BUnit$.}
  By \cref{lem:process-progress}:
  \begin{itemize}
    \item $\E$ a value: $\E = \CUnit$ and $\Proc = \PExp\CUnit$.
    \item $\E \EReduce \E*$: $\Proc \PReduce \PExp{\E*}$ by \textsc{R-Exp}.
    \item $\E = \ECapp{\EApp\C\V}$ blocked: if $\C \in \{\CNew,\CFork\}$ then
      $\Proc$ reduces by \textsc{R-New}, resp.\ \textsc{R-Fork}. Otherwise $\C$
      acts on a channel $\EVar \in \SessCtx$, which no restriction binds, so no
      channel rule applies and $\Blocked\Proc$ by \textsc{B-ExpBlocked}.
  \end{itemize}

  \proofcase{\textsc{TP-Par}, $\Proc = \PPar{\Proc[1]}{\Proc[2]}$.}
  Apply the induction hypothesis to $\Proc[1]$ and to $\Proc[2]$.
  \begin{itemize}
    \item If $\Proc[1] \PReduce \Proc[1]'$, then $\Proc \PReduce
      \PPar{\Proc[1]'}{\Proc[2]}$ by \textsc{R-Par}; symmetrically, if
      $\Proc[2] \PReduce \Proc[2]'$, then $\Proc \PReduce
      \PPar{\Proc[1]}{\Proc[2]'}$ by \textsc{R-Par} together with
      $\Proc \Congruent \PPar{\Proc[2]}{\Proc[1]}$ (\textsc{Par-Comm}) and
      \textsc{R-Struct}.
    \item Otherwise, if $\Proc[1] \Congruent \PExp\CUnit$, then $\Proc \Congruent
      \Proc[2]$ by \textsc{Par-Unit}, and the claim for $\Proc$ is the claim
      already established for $\Proc[2]$; symmetrically if $\Proc[2] \Congruent
      \PExp\CUnit$.
    \item Otherwise $\Proc[1]$ and $\Proc[2]$ are both blocked, and
      $\Blocked\Proc$ by \textsc{B-ParBlocked}.
  \end{itemize}

  \proofcase{\textsc{TP-Res}, $\Proc = \PNew{\BindGroup[1]}{\BindGroup[2]}\ProcQ$.}
  By inversion $\ProcQ$ is typed in $\SessCtx$ extended by $\BindGroup[1],
  \BindGroup[2]$, again a session context. By the induction hypothesis on
  $\ProcQ$:
  \begin{itemize}
    \item $\ProcQ \PReduce \ProcQ'$: then $\Proc \PReduce
      \PNew{\BindGroup[1]}{\BindGroup[2]}{\ProcQ'}$ by \textsc{R-Bind}.
    \item $\ProcQ \Congruent \PExp\CUnit$: excluded. The endpoints of
      $\BindGroup[1],\BindGroup[2]$ carry linear session types $\SSeq\S\STerm$,
      $\SDual{(\SSeq\S\SWait)}$, and $\BindGivesCtx{}{}{}$ requires each to be
      consumed by a process, so a body $\Congruent \PExp\CUnit$ is untypable.
    \item $\ProcQ$ blocked. By the definition of $\Blocked{}$, every process of
      $\ProcQ$ is blocked and every restriction inside $\ProcQ$ is blocked; the
      inner restrictions are closed by the induction hypothesis, so only the
      channel bound \emph{here} remains. This restriction binds one channel with
      two dual endpoints: $\EVar$, the head of $\BindGroup[1]$, governed by $\S$,
      and $\EVary$, the head of $\BindGroup[2]$, governed by $\SDual\S$. (Earlier
      splits may have left further fragments in $\BindGroup[1],\BindGroup[2]$;
      \textsc{R-LSplit} and \textsc{R-RSplit} act on such a fragment in place,
      wherever it sits in its group.) Exactly one of the following holds.
      \begin{enumerate}
        \item A process is blocked on a fragment of $\BindGroup[1]$ or
          $\BindGroup[2]$ with $\CSplit$ or $\CSSplit$, or on $\EVar$ or $\EVary$
          with $\CDrop$ or $\CDiscard$, or on $\EVar$ or $\EVary$ with $\CAcq$
          while the marker $\BLeadSep{}$ is present. The matching rule
          \textsc{R-LSplit}, \textsc{R-RSplit}, \textsc{R-Drop},
          \textsc{R-Discard}, resp.\ \textsc{R-Acquire} fires and $\Proc$ reduces.
        \item Case (1) does not apply, and two distinct processes are blocked at
          communicating actions on the two ends of the bound channel: one on
          $\EVar$ and one on $\EVary$. Since $\EVar$ and $\EVary$ are the two ends
          of a single channel their session types are dual, so 
          the two actions form a matched pair, one of $(\CSend,\CRecv)$,
          $(\CSelect{k},\CBranch)$, $(\CTerm,\CWait)$, and $\Proc$ reduces by
          \textsc{R-Com}, \textsc{R-Choice}, resp.\ \textsc{R-Close}.
        \item Neither (1) nor (2) applies. Then at least one of $\EVar,\EVary$ is
          not blocked at a communicating action, i.e.\ $\EVar \notin \BC(\ProcQ)$
          or $\EVary \notin \BC(\ProcQ)$: the endpoint holds no process ready to
          communicate on it, so no synchronization with the other end is
          possible. This is the premise of \textsc{B-NuBlocked}, which, with
          $\Blocked\ProcQ$, gives $\Blocked\Proc$.
      \end{enumerate}
  \end{itemize}
\end{proof}

\subsection{Soundness and Completeness for Algorithmic Typing}
\label{sec:appendix-metatheory-sc}

%%% Local Variables:
%%% mode: latex
%%% TeX-master: "../popl2027-supplement"
%%% End:
